\subsection{A low-degree class created by the resolved coefficient tail}

\begin{proposition}[A-L11: Proved; the direct low-cohomology comparison fails]
\label{AL11}
Work over $R=\mathbb Z_{(2)}$ and $S=\Gamma_R(T)$.
For every integer $s\geq5$, the construction AT16
with $L=I_1\langle s\rangle$, $m=2$, parameter
$K(2)$, and zero twist can be chosen so that its
torsion-free PD algebra has a nonzero class of
exact order two in form degree one and weight
$2s+7$, already in simplex two.
One may take $X=A=B=G_R$ and both actual target
maps to be the projection killing every new generator.
Thus the resolved-tail algebra does not in general
have the degree-one cohomology of the unresolved
two-generator tail algebra.
The simplex-zero and simplex-one algebras still
equal those of $G_R$; no obstruction to a subsequent
based torsion-$H^1$ saturation is asserted.
The same failure occurs for an actual independent
quadratic class of a free identity-disk partial source,
as constructed at the end of the proof.
\end{proposition}

\begin{proof}
We specify all low-weight generators in the
coefficient resolution used by AT16.  Write
$U=T^{[2]}$, so $T^2=2U$.
Before applying the common weight shift $s$, take
$Q_0$ freely generated over $S$ by $e_{2^j}$,
$j\geq0$, with $e_a\mapsto T^{[a]}$.
This surjects onto $I_1$: the binary expansion of
any positive integer $a$ expresses $T^{[a]}$ as
a product of the corresponding $T^{[2^j]}$
times an odd unit.  The multiplier is odd because
its factorial valuation is
$v_2(a!)-\sum_jv_2((2^j)!)=0$ for the occupied
binary digits.

The only new generators of $Q_1$ in weights at
most four can be chosen as
\[
 \begin{aligned}
 \partial r_2&=Te_1-2e_2,&
 \partial r_3&=Ue_1-Te_2,&
 \partial r_4&=Ue_2-6e_4.
 \end{aligned}                                      \tag{AL.57}
\]
Here and below each subscript is its unshifted
weight.  To verify completeness, the coefficient
augmentation matrices in weights $1,2,3,4$ are
respectively
\[
 (1),\qquad(2,1),\qquad(3,3),\qquad(4,6,1),
 \tag{AL.58}
\]
in the bases
$e_1$;
$Te_1,e_2$;
$Ue_1,Te_2$;
$T^{[3]}e_1,Ue_2,e_4$.
Their kernels are generated in weights two and
three by $r_2,r_3$; in weight four they are generated
by $Ur_2=Tr_3$ and $r_4$.
The coefficient three in the first coordinate of
$Ur_2=(3,-2,0)$ is a unit in $R$, so these last
two vectors generate the entire weight-four kernel.
In particular the generator $e_4$ has been retained.

The only new generators of $Q_2$ up to weight four
and of $Q_3$ up to weight four are consequently
\[
 \partial g_3=Tr_2-2r_3,
 \qquad\partial g_4=Ur_2-Tr_3,
 \qquad\partial h_4=Tg_3-2g_4.                       \tag{AL.59}
\]
Indeed in weight three the two columns of
$Q_1\to Q_0$ are $(2,-2)$ and $(1,-1)$.
In weight four its three columns are
$(3,-2,0),(3,-2,0),(0,1,-6)$, so its kernel
is generated by $(1,-1,0)$; multiplication by
$T$ on $g_3$ is twice this generator.
The two weight-four columns of $Q_2\to Q_1$
are $(2,-2,0)$ and $(1,-1,0)$, giving $h_4$.
This last column is injective, so $Q_4$ has no
generator in weights at most four.
Adding free generators for the remaining kernels
in weights at least five extends these choices
to a free resolution of $I_1$ without changing
any displayed calculation.  Apply the shift $s$
to this entire resolution.

Use $P,W$ and the signs of (AT.75)--(AT.76).
In the bottom module $W^0$, define the classes
\[
 c=[v_{g_3}],\qquad h=[u_{h_4}],\qquad k=[v_{g_4}],
 \qquad a=k-h.
 \tag{AL.60}
\]
The brackets denote the quotient of $P^0$; equivalently
these classes are their images under $d_P$ in $P^1$.
Their weights are $s+3,s+4,s+4$.
The boundary of $v_{h_4}\in P^{-1}$ gives the exact
$S$-linear relation
\[
                    2h+Tc-2k=0,
                 \qquad Tc=2a.                    \tag{AL.61}
\]
There is one $R$-basis vector $c$ in the first
weight, and $h,k$ are an $R$-basis in the next.
For a direct check, the columns of $d_P:P^0\to P^1$
in unshifted weight four, ordered by
$h,Tc,k$, are
\[
 \begin{pmatrix}
 -1&2&0\\
  2&0&2\\
  0&2&1\\
  0&-2&-1\\
  0&0&0
 \end{pmatrix}.                                    \tag{AL.62}
\]
The first and third columns are independent and
have a unit minor, while the middle column is
twice their difference.  There are no bottom
generators in smaller weights.

Consider first the torsion-free image
\[
 E=\operatorname{im}\left(
 \operatorname{FreePD}_{S}(W)\longrightarrow
 \operatorname{FreePD}_{S_{\mathbb Q}}(W_{\mathbb Q})
 \right)
\]
with its untwisted linear differential.  In this
algebra the symbols in (AL.60) are zero-form
generators.  Define
\[
                         \omega=c\,da.
 \tag{AL.63}
\]
Equation (AL.61) gives $2da=Tdc$.  Since odd
squares vanish, $2d\omega=dc\,Tdc=0$.
Torsion-freeness therefore gives $d\omega=0$.
The same relation gives $a\,dc=c\,da$ in the
torsion-free algebra, because twice their difference
is zero.  Consequently
\[
                    d(ca)=2\omega.                 \tag{AL.64}
\]
Its rational primitive is the explicit zero-form
\[
                         F=\frac T2\gamma_2(c),
                         \qquad dF=\omega.         \tag{AL.65}
\]

We show that this primitive is absent integrally.
Give each new generator arity one.  The untwisted
differential and every defining relation preserve
arity.  In arity two and weight $2s+7$, only the
bottom weights $s+3,s+4$ can contribute.  The
degree-zero lattice is therefore generated by
$T\gamma_2(c),ch,ck$.
Inside the rational algebra, (AL.61) gives
\[
                  ck=ch+T\gamma_2(c).
\]
The two remaining vectors are rationally independent:
$c,h$ are independent over $S_{\mathbb Q}$ in this
range, as the $u_{g_4}$ coordinate in (AL.62)
detects $h$ independently of $c$.
Thus the entire lattice in question is exactly
\[
                 R\,ch\ \oplus\ R\,T\gamma_2(c).
 \tag{AL.66}
\]
It does not contain $F$, since $1/2\notin R$.
There are no additional smaller-weight generators;
generators of larger weight cannot enter an arity-two
monomial of this weight.

The rational contraction (AT.80), extended by the
arity derivation, proves that $E_{\mathbb Q}$ has
only coefficient cohomology.  In particular a
zero-form primitive of $\omega$ in arity two is
unique: the difference of two such primitives
would be a closed zero-form of positive arity.
It must equal $F$.  Equation (AL.66) therefore
shows that $\omega$ is not an integral boundary.
Together with (AL.64), this proves its exact order two.

Finally apply AT16 with $X=A=B=G_R$, zero maps
$z,\eta$, and the above resolution.  All new
generator images in the two targets are zero,
so these are actual maps.  Our choice $s\geq5$
puts the first new weight above their cutoffs
through degree four.  In simplex two, $K(2)_2=R$.
The coefficient inclusion $S\to(G_R)_2$ and any
vertex evaluation $(G_R)_2\to S$ give actual
oppositely directed PD cochain maps between $E$
and that simplex algebra of the extension, with
composite the identity on $E$.
They survive the torsion-free quotients.
Thus the class (AL.63) remains nonzero in the
actual AT16 algebra.  The sphere parameter has
zero simplices in degrees zero and one, proving
the stated limitation concerning based saturation.

Here are the actual comparison map and a nonzero
quadratic example.  The resolution augmentation
gives a cochain map
\[
 W\longrightarrow
       \bigl(Lv\xrightarrow{\,2\,}Lu\bigr),
 \qquad
 u_0(q)\longmapsto u(\varepsilon q),\quad
 v_0(q)\longmapsto v(\varepsilon q),                 \tag{AL.67}
\]
and sends every other summand, including $W^0$, to
zero.  The identities $\varepsilon\partial=0$
check every differential.  Its free PD extension
and torsion-free quotient give the indicated
comparison to the unresolved tail algebra; it
sends (AL.63) to zero.  The unresolved zero-twist
algebra has no new generators in forms zero or
one, and hence has zero $H^1$ over $G_R$.

For a nonzero quadratic class, choose $b\geq3$
and set $s=2b$.  Over $G_R$ adjoin the free vertical
identity disk $xK(2)\xrightarrow{d}yK(2)$, with
forms one and two and weight $b$.  Call this
partial source $X$.  In simplex two, let $x_*,y_*$
be the two top normalized generators.
Their faces vanish.  Thus
$\gamma_2(y_*)$ and $x_*y_*$ are normalized
horizontal cycles and define actual simplicial
linear maps from $K(2)$, by the Dold--Kan
correspondence.  Multiplication by coefficients
gives maps on $I_1\langle2b\rangle\otimes K(2)$
with $d\eta=2z$ exactly as in (AT.74).
The two specified maps $X\to G_R$ kill $x,y$.

Apply AT16 with these maps and the same resolution.
The projection killing the old $x,y$ variables
is an actual map to the preceding zero-twist
extension and preserves (AL.63).  The class is
therefore still nonzero.  Formula (AL.67) also
defines an actual map to the unresolved twisted
tail algebra, because its top generator
differentials are now $du=z$ and $dv=2u-\eta$,
the same two twisting equations on both sides.
That algebra has unchanged form-zero and form-one
terms and $H^1(X)=0$, so its $H^1$ is zero.
In simplex two its parameter is precisely the
single independent class $[\gamma_2(y_*)]$ of
A-S09-21.  This gives the explicit failed
degree-one comparison for that quadratic class.
\end{proof}
